\chapter{Soluciones comentadas del banco de test}
\label{ch:test-matmo-2026-sol}
\noindent\textbf{Corrección.} Las letras se refieren al orden a), b), c) del banco. Cada explicación indica el capítulo del manual que debe repasarse.

\paragraph{T01. Respuesta b) -- capítulo 07} 8090 es una aleación Al-Li; la designación no identifica una superaleación de níquel.

\paragraph{T02. Respuesta c) -- capítulo 01} La elección también depende de carga, rigidez, masa, corrosión, vida, fabricación y temperatura real del componente.

\paragraph{T03. Respuesta a) -- capítulo 08} El envejecimiento modifica el estado de precipitados; tiempo y temperatura pueden mejorar o degradar propiedades y no todos los aluminios responden igual.

\paragraph{T04. Respuesta b) -- capítulo 09} Fibras y matriz participan en la respuesta del material. La matriz transmite esfuerzos por cortadura y condiciona la integridad de las interfaces.

\paragraph{T05. Respuesta c) -- capítulo 09} La estructura reticulada afecta al procesado y al comportamiento térmico; no se identifican las familias solo por su rigidez.

\paragraph{T06. Respuesta a) -- capítulo 07} La fase alfa tiene estructura hexagonal compacta y la beta cúbica centrada en el cuerpo. Las fracciones dependen de aleación e historia térmica.

\paragraph{T07. Respuesta b) -- capítulo 09} GLARE es un laminado fibra-metal; no debe asignársele un módulo o comportamiento universal sin indicar orientación y configuración.

\paragraph{T08. Respuesta c) -- capítulo 09} La clasificación diferencia rotura de la interfaz, del material adhesivo y del sustrato.

\paragraph{T09. Respuesta a) -- capítulo 09} Las caras soportan gran parte de la tracción y compresión por flexión, mientras el núcleo mantiene la separación y contribuye a cortante.

\paragraph{T10. Respuesta b) -- capítulo 10} La ausencia de límites de grano ayuda frente a mecanismos determinados de fluencia; no elimina otros mecanismos ni sustituye el diseño térmico.

\paragraph{T11. Respuesta c) -- capítulo 01} La masa mínima idealizada depende de rho dividido por tensión admisible cuando geometría y tipo de carga se mantienen.

\paragraph{T12. Respuesta a) -- capítulo 03} Debe evaluarse $K_I$ frente a la tenacidad del estado de tensión aplicable, comprobando geometría, espesor y validez del modelo.

\paragraph{T13. Respuesta b) -- capítulo 03} Confundir $a$ y $2a$ produce errores importantes en el tamaño crítico y en $K_I$.

\paragraph{T14. Respuesta c) -- capítulo 04} La ley en su región de validez es $\mathrm{d}a/\mathrm{d}N=C(\Delta K)^m$. No debe extrapolarse a umbral o rotura rápida.

\paragraph{T15. Respuesta a) -- capítulo 04} La suma de fracciones de daño es una idealización que no contempla el orden ni todas las interacciones.

\paragraph{T16. Respuesta b) -- capítulo 05} Eliminar bordes de grano no elimina mecanismos intragranulares de deformación dependiente del tiempo.

\paragraph{T17. Respuesta c) -- capítulo 05} Debe declararse la convención del tiempo, el valor de C y el dominio de tensiones y mecanismos.

\paragraph{T18. Respuesta a) -- capítulo 11} Tiempo de vuelo, atenuación, orientación y calibración condicionan la detección y dimensionado.

\paragraph{T19. Respuesta b) -- capítulo 11} Se interpreta la radiación detectada junto a emisividad y condiciones térmicas; una anomalía no identifica por sí sola el defecto.

\paragraph{T20. Respuesta c) -- capítulo 06} La pareja material-medio y la geometría modifican el proceso; los potenciales no dependen solo del nombre de la aleación.

\paragraph{T21. Respuesta a) -- capítulo 21} La aplicabilidad, acciones, revisiones incorporadas y plazos de la directiva deben demostrarse documentalmente.

\paragraph{T22. Respuesta b) -- capítulo 18} Aprobación organizativa y autorización individual de certificación son cuestiones distintas.

\paragraph{T23. Respuesta c) -- capítulo 19} La licencia no concede automáticamente autorización de una organización ni derechos sobre cualquier tipo o tarea.

\paragraph{T24. Respuesta a) -- capítulo 19} La normativa distingue formación aprobada, licencia y aprobación de organización de mantenimiento.

\paragraph{T25. Respuesta b) -- capítulo 16} La gestión de la aeronavegabilidad y la ejecución del mantenimiento son funciones relacionadas, pero distintas.

\paragraph{T26. Respuesta c) -- capítulo 16} ARC acredita el proceso de revisión en su ámbito; CRS documenta la puesta en servicio tras mantenimiento.

\paragraph{T27. Respuesta a) -- capítulo 15} La Parte 21 es el marco para certificación de productos, modificaciones y determinadas aprobaciones y permisos de vuelo.

\paragraph{T28. Respuesta b) -- capítulo 14} 748/2012 organiza reglas de aeronavegabilidad inicial y 1321/2014 de continuada dentro de sus respectivos ámbitos.

\paragraph{T29. Respuesta c) -- capítulo 15} Se diferencia certificado de tipo, certificado individual de aeronavegabilidad y certificación de mantenimiento.

\paragraph{T30. Respuesta a) -- capítulo 14} La masa no determina por sí sola el régimen legal. Es preciso comprobar el artículo de aplicabilidad vigente.

\paragraph{T31. Respuesta b) -- capítulo 17} Deben conservarse todos los contadores, sus referencias y las tolerancias solo si el programa aplicable las contempla.

\paragraph{T32. Respuesta c) -- capítulo 21} La autoridad emisora y la jurisdicción de operación importan; una tabla histórica no demuestra estado legal actual.

\paragraph{T33. Respuesta a) -- capítulo 23} El BOE-A-2022-21487 suministrado corresponde a una convocatoria de personal; no es fundamento de normativa militar.

\paragraph{T34. Respuesta b) -- capítulo 14} Una modificación publicada puede tener distintas fechas de aplicabilidad según artículos y supuestos.

\paragraph{T35. Respuesta c) -- capítulo 19} Deben verificarse categoría y habilitación, condiciones del trabajo y autorización de la organización cuando proceda.

\paragraph{T36. Respuesta a) -- capítulo 22} Una tendencia es evidencia diagnóstica, no una acción de mantenimiento aprobada ni autorización de vida.


\section*{Fuentes abiertas recomendadas para comprobar conceptos}
\begin{itemize}
\item Guía docente oficial: \url{https://www.unileon.es/estudiantes/oferta-academica/grados/grado-en-ingenieria-aeroespacial/plan-estudios?cursoa=2026&id=0710341&te=G}.
\item Mecánica, fractura y fluencia: MIT OpenCourseWare, Roylance: \url{https://ocw.mit.edu/courses/3-11-mechanics-of-materials-fall-1999/pages/modules/}.
\item Ensayos, adhesivos y estructuras: FAA Aviation Maintenance Technician Handbook, \url{https://www.faa.gov/regulations_policies/handbooks_manuals/aviation}.
\item Aeronavegabilidad continuada y aplicabilidad: EASA \url{https://www.easa.europa.eu/en/regulations/continuing-airworthiness}.
\item Certificación de motores: \url{https://www.easa.europa.eu/en/document-library/certification-specifications/group/cs-e-engines}.
\end{itemize}
\noindent\textbf{Advertencia:} confirmar siempre las versiones oficiales aplicables de las normas y las condiciones del caso. Los PDF normativos aportados pueden ser anteriores a la normativa en vigor.
